\documentclass[11pt]{article}
\usepackage[ngerman]{babel}
\usepackage[a4paper,margin=1.1in]{geometry}
\usepackage[T1]{fontenc}
\usepackage[utf8]{inputenc}
\usepackage{lmodern}
\usepackage{amsmath,amssymb,amsthm,mathrsfs}
\usepackage{microtype}
\usepackage{fancyhdr}
\usepackage{array,longtable,enumitem}
\setlength{\parindent}{1.5em}
\setlength{\parskip}{0.18em}
\emergencystretch=3em
\pagestyle{fancy}
\fancyhf{}
\cfoot{\thepage}
\lhead{Heinrich Weber, Lehrbuch der Algebra, Band II}
\rhead{wortgetreue Reparaturabschrift}
\newcommand{\Om}{\Omega}
\newcommand{\Hm}{H_m}
\newcommand{\Lm}{\Omega_m}
\newcommand{\eps}{\varepsilon}
\newcommand{\srcpage}[1]{\par\medskip\noindent\textsf{[Quellseite~#1]}\par\smallskip}
\begin{document}

\section*{\S\,195. Ueber die Classenzahl in dem in $\Omega_m$ enthaltenen reellen Körper.}

\srcpage{711}
Unsere allgemeinen Principien können auch angewandt werden, um die Classenzahlen in den Theilern des Körpers $\Omega_m$ zu bestimmen.

Insbesondere sind die Formeln (6), (7), \S~192 anwendbar, wenn wir die Grade der Primideale in einem solchen Theiler kennen. Die Frage nach dem Verhältniss der Classenzahlen in den Theilern eines Körpers $\Omega_m$ zu der Classenzahl in $\Omega_m$ selbst ist von grossem Interesse und soll hier für den Fall behandelt werden, dass es sich um den in $\Omega_m$ enthaltenen reellen Körper $H_m$ handelt, den wir ja schon im \S~176 f. untersucht haben\footnote{Vgl. Kummer, „Bestimmung der Anzahl u.s.f.“. Crelle's Journ., Bd. 40 (1849).}. Es hat sich dort ergeben, dass $q$ im Körper $H_m$ die $\frac12\mu$te Potenz einer Primzahl ersten Grades ist, dass die anderen Primzahlen $p$ in $e_1$ Primfactoren $f_1$ten Grades zerfallen, so dass
\begin{equation}
f_1e_1=\frac12\mu,
\tag{1}
\end{equation}
und dass zwei Arten von Primzahlen $p_1,p_2$ zu unterscheiden sind, so dass bei den Primzahlen erster Art $f_1=\frac12 f$, $e_1=e$, bei denen der zweiten Art $f_1=f$, $e_1=\frac12 e$ ist.

Für den Körper $H_m$ erhalten wir demnach aus \S~192, (6), (7) die Classenzahl $h_1$ nach der Formel:
\srcpage{712}
\begin{equation}
g_1h_1=\lim_{s=1}(s-1)\Phi_1(s),
\tag{2}
\end{equation}
\begin{equation}
\Phi_1(s)=\frac1{1-q^{-s}}\prod_{p_1}\frac1{(1-p_1^{-\frac12sf})^e}\prod_{p_2}\frac1{(1-p_2^{-sf})^{\frac12e}},
\tag{3}
\end{equation}
worin sich die beiden Productzeichen auf alle Primzahlen $p_1$ und $p_2$ beziehen. Mit Benutzung der oben erklärten Zeichen $f_1,e_1$ kann man dafür auch setzen:
\begin{equation}
\Phi_1(s)=\frac1{1-q^{-s}}\prod_p\frac1{(1-p^{-sf_1})^{e_1}},
\tag{4}
\end{equation}
und das Productzeichen auf alle Primzahlen $p$ erstrecken.

Die Bedeutung von $g_1$ ergiebt sich aus \S~190, (11).

Nun wollen wir aus der Gruppe $C$ der Charaktere $\chi$ einen Theiler $C_1$ aussondern, der aus allen den Charakteren $\chi_1$ bestehen soll, für die
\begin{equation}
\chi_1(-1)=+1
\tag{5}
\end{equation}
ist. Zunächst ist klar, dass diese Charaktere $\chi_1$ eine Gruppe $C_1$ bilden. Diese Gruppe ist aber nicht mit $C$ identisch, weil es gewiss einen Charakter $\chi_0$ giebt, für den $\chi_0(-1)=-1$ ist. Wir brauchen, um einen solchen zu erhalten, nur in \S~193, (6) für $\Theta$ eine primitive $\mu$te Einheitswurzel, oder in \S~193, (8) $\eps=-1$ zu setzen. Dann bildet die Gesammtheit der Charaktere $\chi_0\chi_1$, die alle nicht zu $C_1$ gehören, eine Nebengruppe $C_2$, deren Elemente mit $\chi_2$ bezeichnet werden mögen. Jeder Charakter $\chi$ ist also dann entweder in $C_1$ oder in $C_2$ enthalten, denn wenn $\chi$ nicht in $C_1$ vorkommt, so kommt $\chi_0^{-1}\chi$ darin vor, und folglich $\chi$ in $C_2$. Es ist daher $C_1$ ein Theiler von $C$ vom Index 2.

Die Charaktere $\chi_1$ kann man auf folgende Art bilden:

1) Ist $m$ ungerade, so setzen wir, wie in \S~193,
\[
n\equiv c^\gamma\pmod m,
\]
und lassen in
\begin{equation}
\chi_1(n)=\Theta^\gamma
\tag{6}
\end{equation}
$\Theta$ die Gesammtheit der $\frac12\mu$ten Einheitswurzeln durchlaufen.

2) Ist $m$ gerade, so setzen wir
\[
n\equiv(-1)^\alpha5^\beta,
\]
und erhalten
\begin{equation}
\chi_1(n)=\Theta^\beta,
\tag{7}
\end{equation}
worin $\Theta$ wieder die sämmtlichen $\frac12\mu$ten Einheitswurzeln durchläuft.

\srcpage{713}
Beides ergiebt sich leicht dadurch, dass für $n=-1$ im ersten Falle $\gamma=\frac12\mu$, im zweiten $\alpha=1$, $\beta=0$ ist, also beide Male $\chi_1(-1)=1$ ist, und dass jede der Formeln (6), (7) genau $\frac12\mu$ verschiedene Charaktere darstellt.

Worauf es nun wesentlich ankommt, ist Folgendes:
\begin{center}
\textbf{Ist $p$ irgend eine von $q$ verschiedene Primzahl, so erhält man in der Form $\chi_1(p)$, wenn $\chi_1$ die Gruppe $C_1$ durchläuft, jede $f_1$te Einheitswurzel, und jede gleich oft, also $e_1$ mal.}
\end{center}
Dass alle $\chi_1(p)$ Einheitswurzeln vom Grade $f_1$ sind, ist klar, denn es ist immer (vgl. \S~177) $p^{f_1}\equiv\pm1\pmod m$, und daher
\[
[\chi_1(p)]^{f_1}=\chi_1(\pm1)=1.
\]
Wenn aber unter den $\chi_1(p)$ eine primitive $f_1$te Einheitswurzel vorkommt, so kommen alle anderen $f_1$ten Einheitswurzeln auch darunter vor, und gleich oft, nämlich so oft wie die Einheitswurzel $1$.

Dies folgt unmittelbar aus der Gruppennatur der $\chi_1$. Es ist also nun noch zu zeigen, dass unter den $\chi_1(p)$ eine primitive $f_1$te Einheitswurzel vorkommt.

Diese Möglichkeit ergiebt sich aber aus der Bestimmung von $\chi_1$ sehr einfach. Ist zunächst $\Theta$ in (6) bei ungeradem $m$ eine primitive Einheitswurzel vom Grade $\frac12\mu$, und ist $\gamma$ der Index von $p$, so ist eine Potenz von $\Theta^\gamma$, $\Theta^{\gamma\lambda}$ dann und nur dann gleich $1$, wenn
\[
2\gamma\lambda\equiv0\pmod\mu.
\]
Nun ist $\mu=ef$ und $e$ der grösste gemeinschaftliche Theiler von $\gamma$ und $\mu$; und daher wird diese Bedingung $2\lambda\equiv0\pmod f$. Es ist aber nach \S~177, 3, $f_1=f$ oder $=\frac12f$, je nachdem $f$ ungerade oder gerade ist, und folglich muss $\lambda\equiv0\pmod {f_1}$ sein. Also ist $\Theta^\gamma$ primitive $f_1$te Einheitswurzel.

Ist aber $m$ gerade, so nehmen wir in (7) gleichfalls für $\Theta$ eine primitive $\frac12\mu$te Einheitswurzel und setzen
\[
p\equiv(-1)^\alpha5^\beta\pmod m,
\]
so dass $f$ die kleinste positive Zahl ist, die den Congruenzen
\[
\alpha f\equiv0\pmod2,
\qquad
\beta f\equiv0\pmod{\tfrac12\mu}
\]
genügt. Es ist nun $\Theta^{\beta\lambda}$ nur dann gleich $1$, wenn $\beta\lambda\equiv0\pmod{\frac12\mu}$ ist. Wenn nicht $\lambda=1$, also $\beta\not\equiv0$ ist, so ist der kleinste Werth, den $\lambda$ haben kann, gerade, und daher ist dieser kleinste Werth von $\lambda=f$.

\srcpage{714}
Wenn aber $\lambda=1$ ist, so muss $\beta=0$ sein, und es sind noch zwei Fälle möglich:
\[
\begin{array}{ll}
entweder & f=1,\quad \alpha=0,\quad p\equiv1\pmod m,\quad \lambda=f,\\[2mm]
oder & f=2,\quad \alpha=1,\quad p\equiv-1\pmod m,\quad \lambda=\frac12f.
\end{array}
\]
Nur in diesem letzten Falle gehört $p$ zur ersten Art, und es folgt also, dass auch hier allgemein $\Theta^\beta$ eine primitive $f_1$te Einheitswurzel ist. Hiermit ist dann der Satz bewiesen.

Dieser Satz gestattet uns die Anwendung der Formel \S~193, (9), die hier ergiebt:
\begin{equation}
(1-p^{-sf_1})^{e_1}=\prod_{\chi_1}[1-\chi_1(p)p^{-s}],
\tag{8}
\end{equation}
worin sich das Product auf alle Charaktere $\chi_1$ der Gruppe $C_1$ erstreckt, und nun lässt sich die Function $\Phi_1(s)$ ebenso wie $\Phi(s)$ im \S~193 weiter entwickeln. Man erhält auf demselben Wege, wie dort, die Formel (11),
\begin{equation}
\Phi_1(s)=\frac1{1-q^{-s}}\prod_{\chi_1}\sum_{n=1}^{\infty}\frac{\chi_1(n)}{n^s},
\tag{9}
\end{equation}
und wenn man zur Grenze $s=1$ übergeht,
\begin{equation}
g_1h_1=\prod_{\chi_1}\sum_{n=1}^{\infty}\frac{\chi_1(n)}{n},
\tag{10}
\end{equation}
worin aber jetzt das Product auf alle Charaktere der Gruppe $C_1$ mit Ausnahme des Hauptcharakters zu erstrecken ist. Dieses Product ist also ein Theil des Productes, durch welches im \S~193 die Zahl $gh$ ausgedrückt ist.

Die Zahlen $g,g_1$ ergeben sich aus dem im \S~190, (11) angegebenen allgemeinen Ausdrucke:
\begin{equation}
\frac{2^\nu\pi^{n-\nu}L}{w\sqrt{\pm\Delta}}.
\tag{11}
\end{equation}
Da, wie im \S~178 nachgewiesen ist, die Einheiten in $\Omega_m$ und $H_m$, von den Einheitswurzeln abgesehen, dieselben sind, so ist auch ein Fundamentalsystem von Einheiten in $H_m$ zugleich ein Fundamentalsystem in $\Omega_m$.

Bedeutet $\eps_1,\eps_2,\ldots,\eps_{\nu-1}$ daher ein Fundamentalsystem reeller Einheiten, und sind $\eps_{i,1},\eps_{i,2},\ldots,\eps_{i,\nu}$ die conjugirten Einheiten zu $\eps_i$, so sind die conjugirten Logarithmen im Körper $H_m$:
\[
\log|\eps_{i,1}|,\quad\log|\eps_{i,2}|,\quad\ldots,\quad\log|\eps_{i,\nu}|,
\]
dagegen im Körper $\Omega_m$, in dem nur conjugirt imaginäre Paare vorkommen (\S~185),
\[
2\log|\eps_{i,1}|,\quad2\log|\eps_{i,2}|,\quad\ldots,\quad2\log|\eps_{i,\nu}|.
\]

\srcpage{715}
Bezeichnen wir daher mit
\begin{equation}
E=\pm\begin{vmatrix}
\log|\eps_{1,1}|&\log|\eps_{1,2}|&\cdots&\log|\eps_{1,\nu-1}|\\
\log|\eps_{2,1}|&\log|\eps_{2,2}|&\cdots&\log|\eps_{2,\nu-1}|\\
\cdots&\cdots&\cdots&\cdots\\
\log|\eps_{\nu-1,1}|&\log|\eps_{\nu-1,2}|&\cdots&\log|\eps_{\nu-1,\nu-1}|
\end{vmatrix}
\tag{12}
\end{equation}
den Regulator des Körpers $H_m$ (\S~187), so ist bei der Anwendung des Ausdruckes (11) im Körper $\Omega_m$
\[
L=2^{\nu-1}E,
\]
und im Körper $H_m$
\[
L=E
\]
zu setzen. Es ist ferner
\[
\begin{array}{lll}
in\ \Omega_m: & n=\mu, & \nu=\frac12\mu,\\
in\ H_m: & n=\frac12\mu, & \nu=\frac12\mu.
\end{array}
\]
Im Körper $H_m$ sind nur die zwei Einheitswurzeln $\pm1$ enthalten. In $\Omega_m$ sind die mit positivem und negativem Zeichen genommenen Potenzen von $r$, sonst aber keine Einheitswurzeln enthalten (\S~175), und bei der Zählung ist noch zu beachten, dass bei geradem $m$ unter den Potenzen von $r$ auch $-1$ enthalten ist, bei ungeradem $m$ nicht. Daher haben wir
\[
\begin{array}{lll}
in\ \Omega_m: & w=2m, & m\text{ ungerade},\\
& w=m, & m\text{ gerade},\\
in\ H_m: & w=2.&
\end{array}
\]
Endlich haben wir noch, wenn $\Delta,\Delta_1$ die Grundzahlen der Körper $\Omega_m,H_m$ sind (\S~176):
\[
\pm\Delta=q\Delta_1^2\quad m\text{ ungerade},
\qquad
\pm\Delta=4\Delta_1^2\quad m\text{ gerade},
\]
und es folgt, wenn wir der Einfachheit wegen $\nu$ für $\frac12\mu$ setzen:
\begin{align}
g&=\frac{2^{2(\nu-1)}\pi^\nu E}{m\sqrt q\,\Delta_1}, &&m\text{ ungerade},\notag\\
g&=\frac{2^{2(\nu-1)}\pi^\nu E}{m\Delta_1}, &&m\text{ gerade},\notag\\
g_1&=\frac{2^{\nu-1}E}{\sqrt{\Delta_1}},\tag{13}
\end{align}
woraus
\begin{align}
g&=\frac{2^{\nu-1}\pi^\nu g_1}{m\sqrt{q\Delta_1}}, &&m\text{ ungerade},\notag\\
g&=\frac{2^{\nu-1}\pi^\nu g_1}{m\sqrt{\Delta_1}}, &&m\text{ gerade}.\tag{14}
\end{align}

\srcpage{716}
Wenn wir die Formel (10) mit der Formel \S~193, (15) verbinden und die im \S~194 eingeführte Bezeichnung
\[
X_1=\sum\frac{\chi_1(n)}{n},
\qquad
X_2=\sum\frac{\chi_2(n)}{n}
\]
gebrauchen, so ergiebt sich
\[
gh=g_1h_1\prod X_2.
\]
Ersetzt man $g$ durch seinen Werth (14), so folgt für die Classenzahl
\begin{equation}
h=kh_1,
\tag{15}
\end{equation}
wenn
\begin{equation}
\begin{aligned}
k&=\frac{m\sqrt{q\Delta_1}}{2^{\nu-1}\pi^\nu}\prod X_2 &&\text{bei ungeradem }m,\\
&=\frac{m\sqrt{\Delta_1}}{2^{\nu-1}\pi^\nu}\prod X_2 &&\text{bei geradem }m,
\end{aligned}
\tag{16}
\end{equation}
und worin [nach (10), (13)]:
\begin{equation}
h_1=\frac{\sqrt{\Delta_1}}{2^{\nu-1}E}\prod X_1
\tag{17}
\end{equation}
die Classenzahl des Körpers $H_m$ ist. Im letzteren Ausdrucke erstreckt sich das Product auf alle Charaktere $\chi_1$ der Gruppe $C_1$ mit Ausnahme des Hauptcharakters.

Dass $k$ eine ganze Zahl ist, und daher $h$ ein Vielfaches von $h_1$, wird sich bald zeigen. Die Zahlen $k$ und $h_1$ heissen der erste und zweite Factor der Classenzahl.

In dem Falle, dass
\[
m=2^\varkappa
\]
eine Potenz von 2 ist, haben wir nach \S~176, (7):
\[
\Delta_1=2^{(\varkappa-1)2^{\varkappa-2}-1},
\]
und wenn wir diesen Werth einsetzen, so können wir für die vorstehenden Formeln in diesem Falle schreiben:
\begin{equation}
\begin{aligned}
k&=\frac{\sqrt2\,2^{(\varkappa-3)2^{\varkappa-3}}2^\varkappa}{\pi^\nu}\prod X_2,\\
h_1&=\frac{\sqrt2\,2^{(\varkappa-3)2^{\varkappa-3}}}{E}\prod X_1.
\end{aligned}
\tag{18}
\end{equation}

\section*{\S\,196. Classenzahl im Körper der achten Einheitswurzeln.}

\srcpage{717}
Wir wollen, einerseits zur Veranschaulichung der bisherigen Resultate, andererseits um für die später anzuwendende vollständige Induction eine Basis zu gewinnen, den Fall $m=8$, $\mu=4$, $\nu=2$ zunächst behandeln.

Hier haben wir ausser dem Hauptcharakter nur drei Charaktere, nämlich, wenn
\[
n\equiv(-1)^\alpha5^\beta\pmod8
\]
ist,
\[
\chi_1(n)=(-1)^\beta,
\qquad
\chi_2(n)=(-1)^\alpha,
\qquad
\chi_3(n)=(-1)^{\alpha+\beta}.
\]
Für $n=-1$ ist $\alpha=1$, $\beta=0$, und folglich gehört $\chi_1$ zur ersten, $\chi_2,\chi_3$ zur zweiten Art [\S~195, (5)], und es ergeben sich für die vier Zahlen $n=1,3,5,7$ folgende Werthe dieser Charaktere:
\[
\begin{array}{c|rrrr}
&1&3&5&7\\\hline
\chi_1&+1&-1&-1&+1\\
\chi_2&+1&-1&+1&-1\\
\chi_3&+1&+1&-1&-1.
\end{array}
\]
Daraus erhält man die drei Functionen $f_1,f_2,f_3$ [\S~194, (2)]:
\[
\begin{aligned}
f_1(x)&=x-x^3-x^5+x^7=x(1-x^2)(1-x^4),\\
f_2(x)&=x-x^3+x^5-x^7=x(1-x^2)(1+x^4),\\
f_3(x)&=x+x^3-x^5-x^7=x(1+x^2)(1-x^4).
\end{aligned}
\]
Für $r$ können wir setzen:
\[
r=\frac{1+i}{\sqrt2},\quad r^2=i,
\quad r^3=-\frac{1-i}{\sqrt2},\quad r^4=-1,
\]
und danach ergiebt sich:
\[
\begin{array}{lll}
f_1(r)=2\sqrt2,& f_2(r)=0,& f_3(r)=2\sqrt2\,i,\\
f_1(r^2)=0,& f_2(r^2)=4i,& f_3(r^2)=0,\\
f_1(r^3)=-2\sqrt2,& f_2(r^3)=0,& f_3(r^3)=2\sqrt2\,i.
\end{array}
\]
Für $x=r^4$ verschwinden alle drei Functionen $f(x)$, und für $x=r^5,r^6,r^7$ erhält $f_1(x)$ dieselben, $f_2(x)$ und $f_3(x)$ die entgegengesetzten Werthe, wie für $x=r^3,r^2,r$.

\srcpage{718}
Demnach ergiebt sich aus \S~194, (12) mit Rücksicht auf die trigonometrische Formel
\[
\sin\frac{3\pi}{8}=\cos\frac\pi8
\]
\[
X_1=-\frac1{\sqrt2}\log\tan\frac\pi8,
\qquad
X_2=-\frac\pi4,
\qquad
X_3=-\frac\pi{2\sqrt2}.
\]
Nun ist
\[
\tau=\frac{r-1}{r^2(r+1)}=\tan\frac\pi8=\sqrt2-1
\]
eine Einheit des Körpers $H_8$, der hier nichts Anderes ist, als der aus $\sqrt2$ entspringende quadratische Körper, und die Einheiten darin erhält man also aus den Lösungen der Pell'schen Gleichung
\[
u^2-2v^2=\pm1
\]
in der Form $u+v\sqrt2$.

Die kleinste positive Lösung dieser Pell'schen Gleichung ist aber $u=1$, $v=1$, und folglich findet man nach Bd. I, \S~128 alle Einheiten in $H_8$ in der Form $\pm(1+\sqrt2)^a$, worin $a$ ein positiver oder negativer ganzzahliger Exponent ist; $\tau$ ergiebt sich hieraus für $a=-1$, und daher sind alle Einheiten auch in der Form
\[
\pm\tau^a
\]
enthalten. Es ist also $\tau$ eine fundamentale Einheit, und die in den Formeln des vorigen Paragraphen vorkommende Determinante $E$ [\S~195, (12)] wird hier (da $\tau$ ein echter Bruch ist):
\[
E=-\log\tau.
\]
Man erhält also aus \S~195, (15), (18):
\[
h_1=1,\quad k=1,\quad h=1.
\]
Der Körper $\Omega_8$ ist also ein einclassiger, d. h. ein solcher, in dem die Zerlegung der ganzen Zahlen in Primfactoren nach denselben Gesetzen, wie im Körper der rationalen Zahlen, möglich ist.

Es hat sich dabei zugleich ergeben, dass alle Einheiten im Körper $H_8$ in der Form $\pm\tau^a$ darstellbar sind. Die conjugirten Werthe der Einheit $\tau$ sind aber
\[
\tau_1=\tan\frac\pi8,
\qquad
\tau_3=-\tan\frac{3\pi}{8}=-\tau_1^{-1},
\]
und haben also verschiedene Zeichen. Wir schliessen daraus für diesen Fall auf den Satz:
\begin{center}
\textbf{Eine Einheit in $H_8$, die mit ihren Conjugirten von einerlei Zeichen ist, ist, vom Zeichen abgesehen, das Quadrat einer Einheit.}
\end{center}

\section*{\S\,197. Recurrente Berechnung der Classenzahl im Körper $\Omega_m$, wenn $m$ eine Potenz von 2 ist.}

\srcpage{719}
Zur allgemeinen Bestimmung der Classenzahl im Körper $\Omega_m$ unter der Voraussetzung, dass $m$ irgend eine Potenz von 2 und grösser als 8 ist, wenden wir ein recurrirendes Verfahren an. Es sei also
\begin{equation}
m=2^\varkappa,
\quad \mu=2^{\varkappa-1},
\quad \nu=2^{\varkappa-2},
\quad \nu'=2^{\varkappa-3}.
\tag{1}
\end{equation}
Neben dem Körper $\Omega_m$ betrachten wir den Körper $\Omega_\mu$, der ein Theiler von $\Omega_m$ ist, und den wir hier auch mit $\Omega'_m$ bezeichnen wollen. Es sei $h'$ die Classenzahl im Körper $\Omega'_m$, d. h. die Zahl, die sich aus $h$ ergiebt, wenn $\varkappa$ in $\varkappa-1$ verwandelt wird, und wir setzen
\begin{equation}
h=h'H.
\tag{2}
\end{equation}
Setzen wir $h'$ als schon bekannt voraus, so kommt es nur noch auf die Berechnung von $H$ an, was, wie sich zeigen wird, eine ganze Zahl ist. Indem wir $h$ und $h'$ wie im \S~195 zerlegen, setzen wir
\begin{equation}
h=kh_1,
\qquad
h'=k'h'_1,
\qquad
H=AB,
\tag{3}
\end{equation}
\begin{equation}
k=k'A,
\qquad
h_1=h'_1B,
\tag{4}
\end{equation}
worin $k'$ aus $k$ und $h'_1$ aus $h_1$ durch die Vertauschung von $\varkappa$ mit $\varkappa-1$ hervorgeht. Wir wollen überhaupt jetzt durch Accente an den Buchstaben andeuten, dass $\varkappa-1$ statt $\varkappa$ gesetzt sein soll.

Nach \S~176, (7) ist dann
\[
\Delta_1=2^{(\varkappa-1)2^{\varkappa-2}-1},
\qquad
\Delta'_1=2^{(\varkappa-2)2^{\varkappa-3}-1},
\]
und folglich
\begin{equation}
\Delta_1=2^{\varkappa2^{\varkappa-3}}\Delta'_1.
\tag{5}
\end{equation}
Hat $E$ die Bedeutung \S~195, (12), und geht $E'$ aus $E$ hervor, wenn $\varkappa$ durch $\varkappa-1$ ersetzt wird, so setzen wir noch
\begin{equation}
E=E'D,
\tag{6}
\end{equation}

\srcpage{720}
wodurch $D$ als eine von Null verschiedene positive Zahl definirt ist. Was die Summen
\[
X=\sum\frac{\chi(n)}{n}
\]
betrifft, so ist daran zu erinnern, dass $\chi(n)$ [nach \S~193, (8)] definirt ist durch
\begin{equation}
n\equiv(-1)^\alpha5^\beta\pmod m,
\qquad
\chi(n)=(\pm1)^\alpha\Theta^\beta,
\tag{7}
\end{equation}
worin $\Theta$ eine $\nu$te Einheitswurzel bedeutet.

Unter den $\nu$ten Einheitswurzeln sind aber auch die $\nu'$ten Einheitswurzeln enthalten, und unter den Charakteren $\chi$ auch die Charaktere für den Körper $\Omega'_m$. Wir unterscheiden demnach primitive und imprimitive Charaktere des Körpers $\Omega_m$, indem wir die dem Körper $\Omega_m$ eigenthümlichen Charaktere $\chi$, d. h. die, in denen $\Theta$ eine primitive $\nu$te Einheitswurzel ist, als primitiv bezeichnen.

Unter den Summen $X$ kommen auch alle zu dem Körper $\Omega'_m$ gehörigen Summen vor:
\[
X'=\sum\frac{\chi'(n)}{n}.
\]
Wenn man nun in (4) für $k,k'$ und $h_1,h'_1$ die im \S~195, (16), (17) gebildeten Ausdrücke einsetzt, und dann die Summen $X'$ weghebt, so ergiebt sich:
\begin{equation}
\begin{aligned}
\pi^{\nu'}A&=2\cdot2^{2^{\varkappa-4}(\varkappa-2)}\prod X_2,\\
DB&=2^{2^{\varkappa-4}(\varkappa-2)}\prod X_1,
\end{aligned}
\tag{8}
\end{equation}
worin sich jetzt aber die Producte $\prod$ nur noch auf die mit den primitiven Charakteren $\chi_1,\chi_2$ gebildeten Summen $X_1,X_2$ erstrecken.

In den im \S~194, (12) gegebenen Ausdrücken für die $X$:
\begin{equation}
\begin{aligned}
X_1&=-\frac1{2m}\sum_{s=1}^{m-1} f_1(r^s)\log\left(\sin\frac{s\pi}{m}\right)^2,\\
X_2&=-\frac{i\pi}{m^2}\sum_{s=1}^{m-1}s f_2(r^s)
\end{aligned}
\tag{9}
\end{equation}
treten nun, unter der Voraussetzung, dass $m$ eine Potenz von 2 und dass $\chi_1,\chi_2$ primitive Charaktere sind, bedeutende Vereinfachungen ein.

\srcpage{721}
Die Zahl $1+\mu$ hat die Indices $\alpha=0$, $\beta=\nu'$, und folglich ist nach (7)
\[
\chi(1+\mu)=\Theta^{\nu'}=-1,
\]
wenn $\chi$ ein primitiver Charakter ist. Ferner ist für ein ungerades $n$:
\[
n(1+\mu)\equiv n+\mu\pmod m,
\]
und folglich
\[
\chi(n+\mu)=-\chi(n).
\]
Wenn wir nun die Function $f(r^s)$ betrachten:
\begin{equation}
f(r^s)=\sum_n\chi(n)r^{sn},
\tag{10}
\end{equation}
worin $n$ die ungeraden Zahlen eines vollen Restsystems für den Modul $m$ durchläuft, so erhalten wir, da wir in (10) unter dem Summenzeichen $n$ durch $n+\mu$ ersetzen dürfen:
\[
f(r^s)=-r^{\mu s}\sum\chi(n)r^{sn},
\]
und da $r^\mu=-1$ ist:
\[
f(r^s)=-(-1)^sf(r^s).
\]
Es ist also
\begin{equation}
f(r^s)=0,\quad \text{wenn }s\text{ gerade ist,}
\tag{11}
\end{equation}
und demnach können wir in den Formeln (9) den Summationsbuchstaben $s$ auf die ungeraden Zahlen beschränken, die kleiner als $m$ sind.

Es ist aber nach (10):
\[
f(r^s)=\chi(s)^{-1}\sum_n\chi(ns)r^{sn},
\]
und wenn man $sn$ durch $n$ ersetzt, was bei ungeradem $s$ gestattet ist:
\begin{equation}
f(r^s)=\chi(s)^{-1}f(r),\quad \text{wenn }s\text{ ungerade ist.}
\tag{12}
\end{equation}
Die Function $f(r)$ hat, wenn $\alpha,\beta$ die Indices von $n$ sind, den Ausdruck
\begin{equation}
f(r)=\sum(\pm1)^\alpha\Theta^\beta r^n,
\tag{13}
\end{equation}
und ist also mit der im \S~17 dieses Bandes betrachteten Kreistheilungsresolvente
\[
(\pm1,\Theta,r)
\]
gleichbedeutend, für die im \S~17, (17) die Relation aufgestellt war:
\begin{equation}
(\pm1,\Theta,r)(\pm1,\Theta^{-1},r)=\pm m.
\tag{14}
\end{equation}
Hier gelten durchweg die oberen Zeichen für die Summen $X_1$, die unteren für die Summen $X_2$ [\S~195, (7)].

Wir theilen jetzt die Reihe der in (9) vorkommenden Zahlen $s$ in vier Theile. Es soll $t$ ein Zeichen sein, welches die
\srcpage{722}
Reihe der positiven ungeraden Zahlen von 1 bis $\nu-1$ durchläuft; dann durchläuft das ungerade $s$ die vier Zahlenreihen:
\[
t,
\quad t+\mu,
\quad m-t,
\quad m-t-\mu,
\]
und es ist:
\[
\begin{gathered}
\chi(t+\mu)=-\chi(t),\\
\chi_1(m-t)=\chi_1(t),\qquad \chi_1(m-\mu-t)=-\chi_1(t),\\
\chi_2(m-t)=-\chi_2(t),\qquad \chi_2(m-\mu-t)=\chi_2(t).
\end{gathered}
\]
Nach (12) ergiebt sich hieraus:
\[
\begin{gathered}
f(r^{t+\mu})=-f(r^t),\\
f_1(r^{m-t})=f_1(r^t),\qquad f_1(r^{m-\mu-t})=-f_1(r^t),\\
f_2(r^{m-t})=-f_2(r^t),\qquad f_2(r^{m-\mu-t})=f_2(r^t).
\end{gathered}
\]
Theilt man demnach die Summen (9) in je vier Theile und benutzt noch die trigonometrische Gleichung
\[
\sin\frac{(t+\mu)\pi}{m}=\cos\frac{t\pi}{m},
\]
so ergiebt sich
\[
X_1=-\frac1m\sum_t f_1(r^t)\log\left(\tan\frac{t\pi}{m}\right)^2,
\qquad
X_2=\frac{i\pi}{m}\sum_t f_2(r^t).
\]
In dieser Summe ist
\[
\frac{t\pi}{m}<\frac\pi4,
\qquad
\tan\frac{t\pi}{m}>0,
\]
und folglich erhalten wir mit Rücksicht auf (12), (13), wenn wir $\Theta^{-1}$ für $\Theta$ setzen, wodurch $\chi(t)^{-1}$ in $\chi(t)$ übergeht:
\begin{equation}
X_1=-\frac2m(+1,\Theta^{-1},r)
\sum_{1\le t\le\nu-1}\chi_1(t)\log\tan\frac{t\pi}{m},
\tag{15}
\end{equation}
\begin{equation}
X_2=\frac{i\pi}{m}(-1,\Theta^{-1},r)
\sum_{1\le t\le\nu-1}\chi_2(t).
\tag{16}
\end{equation}

\section*{\S\,198. Der Classenzahlfactor $A$.}

Um nun zunächst $A$ zu berechnen, haben wir den Ausdruck (16) für $X_2$ in die erste Formel (8) des vorigen Paragraphen einzusetzen. Diese Formel können wir auch so darstellen:
\begin{equation}
\pi^{\nu'}A=2m^{1/4}2^{-\nu'}\prod X_2.
\tag{1}
\end{equation}

\srcpage{723}
Bezeichnen wir mit $\alpha,\beta$ die Indices von $t$, so ist $\chi_2(t)=(-1)^\alpha\Theta^\beta$;
\begin{equation}
\sum_t\chi_2(t)=\sum(-1)^\alpha\Theta^\beta=\varphi(\Theta)
\tag{2}
\end{equation}
ist eine ganze Zahl des Körpers $\Omega_\nu$, und es wird:
\[
X_2=\frac{i\pi}{m}(-1,\Theta^{-1},r)\varphi(\Theta).
\]
Nun hat man für $\Theta$ alle Wurzeln der Gleichung
\begin{equation}
\Theta^{\nu'}+1=0
\tag{3}
\end{equation}
zu setzen, und erhält für $A$, mit Rücksicht auf \S~197, (14):
\begin{equation}
A=2^{1-\nu'}\prod_\Theta\varphi(\Theta).
\tag{4}
\end{equation}
Die Summe, durch die in (2) die Zahl $\varphi(\Theta)$ definirt ist, enthält $\nu'$ Glieder von der Form $(-1)^\alpha\Theta^\beta$, worin $\alpha,\beta$ so zu bestimmen sind, dass
\begin{equation}
(-1)^\alpha5^\beta\equiv t\pmod m
\tag{5}
\end{equation}
und $t$ zwischen 0 und $\nu$ liegt. Nun ist zunächst ersichtlich, dass unter den Exponenten $\beta$ nicht zwei nach dem Modul $\nu'$ congruente vorkommen können. Denn ist $\beta\equiv\beta'\pmod{\nu'}$, so ist $5^\beta\equiv5^{\beta'}\pmod\mu$, und aus
\[
(-1)^\alpha5^\beta\equiv t,
\qquad
(-1)^{\alpha'}5^{\beta'}\equiv t'\pmod m
\]
folgt:
\[
(-1)^\alpha t-(-1)^{\alpha'}t'\equiv0\pmod\mu.
\]
Da aber $t$ und $t'$ absolut kleiner als $\nu=\frac12\mu$ sind, so ist dies nur möglich, wenn $t=t'$, $\alpha=\alpha'$ ist. Demnach durchläuft $\beta$ in (2) ein volles Restsystem nach dem Modul $\nu'$.

Da aber $\beta$ nach dem Modul $\nu$ zu nehmen ist, so werden die nach (5) bestimmten $\beta$ theils kleiner, theils grösser als $\nu'$ ausfallen, und wenn wir also $\beta_1$ die Reihe der Zahlen $0,1,\ldots,\nu'-1$ durchlaufen lassen, so erhalten wir in (2) zweierlei Arten von Gliedern:
\begin{equation}
\begin{array}{ll}
1. & (-1)^\alpha\Theta^\beta=(-1)^\alpha\Theta^{\beta_1},\quad \beta=\beta_1<\nu',\\[1mm]
2. & (-1)^\alpha\Theta^\beta=-(-1)^\alpha\Theta^{\beta_1},\quad \beta=\beta_1+\nu'\geqq\nu'.
\end{array}
\tag{6}
\end{equation}
Nun ist nach (5) der absolut kleinste Rest von $(-1)^\alpha5^\beta$ nach dem Modul $m$ positiv, und wenn wir daher eine Zahl $c_\beta=\pm1$ so bestimmen, dass der absolut kleinste Rest von $c_\beta5^\beta$ positiv wird, so ist in dem Falle (6), 1,
\[
c_{\beta_1}=(-1)^\alpha.
\]

\srcpage{724}
Ferner ist $5^{\nu'}\equiv1\pmod\mu$, und da $5^{\nu'}$ nicht auch nach dem Modul $m$ mit 1 congruent ist, so folgt:
\begin{equation}
5^{\nu'}\equiv1+\mu\pmod m.
\tag{7}
\end{equation}
Demnach ist im Falle (6), 2,
\[
(-1)^\alpha5^\beta=(-1)^\alpha5^{\beta_1}5^{\nu'}\equiv(-1)^\alpha5^{\beta_1}+\mu\pmod m,
\]
und folglich ist hier der absolut kleinste Rest von $(-1)^\alpha5^{\beta_1}$ nach dem Modul $m$ negativ. Es ist also im Falle (6), 2,
\[
c_{\beta_1}=-(-1)^\alpha
\]
zu setzen.

Daraus ergiebt sich nach (2) (wenn wir wieder $\beta$ für $\beta_1$ schreiben):
\begin{equation}
\varphi(\Theta)=\sum_{\beta=0}^{\nu'-1}c_\beta\Theta^\beta
=c_0+c_1\Theta+c_2\Theta^2+\cdots+c_{\nu'-1}\Theta^{\nu'-1},
\tag{8}
\end{equation}
und hierin ist
\begin{equation}
c_\beta=+1\quad\text{oder}\quad=-1,
\tag{9}
\end{equation}
je nachdem der absolut kleinste Rest von $5^\beta$ positiv oder negativ ist. Hiernach ist also $\varphi(\Theta)$ eine ganze Zahl des Körpers $\Omega_\nu$.

Es kommt noch darauf an, das Verhalten dieser Zahl zu der Zahl 2, oder vielmehr zu dem in $\Omega_\nu$ enthaltenen Primfactor $(1-\Theta)$ von 2 zu ermitteln. Bilden wir zu diesem Zwecke
\begin{equation}
(1-\Theta)\varphi(\Theta)=2\psi(\Theta),
\tag{10}
\end{equation}
so ergiebt sich
\begin{equation}
\psi(\Theta)=\frac{c_0+c_{\nu'-1}}2+\frac{c_1-c_0}2\Theta+\frac{c_2-c_1}2\Theta^2+\cdots+
\frac{c_{\nu'-1}-c_{\nu'-2}}2\Theta^{\nu'-1}.
\tag{11}
\end{equation}
Die Coefficienten in diesem Ausdrucke sind alle $=0$ oder $=1$, und daher ist auch $\psi(\Theta)$ eine ganze Zahl des Körpers $\Omega_\nu$. Für diese Zahl ergiebt sich aber, wenn man $\Theta\equiv1$ setzt:
\begin{equation}
\psi(\Theta)\equiv c_{\nu'-1}\equiv\pm1\pmod{(1-\Theta)},
\tag{12}
\end{equation}
und folglich ist $\psi(\Theta)$ relativ prim zu 2.

Das in (4) vorkommende Product ist nichts Anderes, als die in Bezug auf den Körper $\Omega_\nu$ genommene Norm der Zahl $\varphi(\Theta)$, und wenn wir diese Norm mit $N$ bezeichnen, so ist nach \S~169
\[
N(1-\Theta)=2,
\]

\srcpage{725}
und folglich nach (10)
\[
N\varphi(\Theta)=2^{\nu'-1}N\psi(\Theta).
\]
Daraus ergiebt sich
\begin{equation}
A=N\psi(\Theta),
\tag{13}
\end{equation}
und hieraus folgt mit Rücksicht auf (12), dass $A$ eine ungerade ganze Zahl ist.

Bezeichnen wir die zu $m=2^\varkappa$ gehörige Zahl $A$ mit $A_\varkappa$, so ergiebt sich aus \S~197, (4) durch vollständige Induction der Ausdruck für den ersten Factor der Classenzahl:
\begin{equation}
k=A_4A_5\cdots A_\varkappa,
\tag{14}
\end{equation}
und daraus der Satz:
\begin{center}
\textbf{1. Der erste Factor der Classenzahl ist eine ungerade ganze Zahl.}
\end{center}
Die Zahl $A_\varkappa$ ist nach der Formel (10) für die ersten Werthe von $\varkappa$ nicht schwer zu berechnen.

Für $m=16$ findet man $\nu'=2$, $\beta=0,1$, $c_0=c_1=1$, folglich
\[
\psi(\Theta)=1,
\qquad
A_4=1.
\]
Für $m=32$, $\nu'=4$, $\beta=0,1,2,3$,
\[
c_0=1,
\quad c_1=1,
\quad c_2=-1,
\quad c_3=-1,
\]
\[
\psi(\Theta)=-\Theta^2,
\]
also auch $A_5=1$.

Für $m=64$, $\nu'=8$, $\beta=0,1,2,3,4,5,6,7$ sind die absolut kleinsten Reste von $5^\beta$
\[
1,\;5,\;25,\;-3,\;-15,\;-11,\;9,\;-19;
\]
also sind die $c_\beta$:
\[
+1,+1,+1,-1,-1,-1,+1,-1
\]
und
\[
\psi(\Theta)=-\Theta^3+\Theta^6-\Theta^7,
\qquad
\psi(\Theta)\psi(-\Theta)=\Theta^6(\Theta^6-2i),
\]
woraus sich leicht ergiebt:
\[
A_6=17.
\]
Eine etwas längere Rechnung ergiebt noch
\[
A_7=21121.
\]

\section*{\S\,199. Der Classenzahlfactor $B$.}

\srcpage{726}
Auf ganz andere Weise muss der Classenzahlfactor $B$ berechnet werden. Hier ist, wenn $\Theta$ wieder eine Primitivwurzel der Gleichung \S~198, (3) ist ($\Theta^{\nu'}+1=0$), und
\begin{equation}
t\equiv(-1)^\alpha5^\beta\pmod m,
\tag{1}
\end{equation}
\begin{equation}
\chi_1(t)=\Theta^\beta
\tag{2}
\end{equation}
zu setzen, und die Formel \S~197, (15) ergiebt:
\begin{equation}
X_1=-\frac2m(1,\Theta^{-1},r)\sum_t\Theta^\beta\log\tan\frac{t\pi}{m}.
\tag{3}
\end{equation}
Nun ist für jedes ungerade $t$, wenn [mit Rücksicht auf (1)]
\[
r=e^{2\pi i/m},
\qquad
r^\nu=i,
\qquad
r^{t\nu}=(-1)^\alpha i
\]
gesetzt wird,
\begin{equation}
\tan\frac{t\pi}{m}=(-1)^\alpha r^t\frac{1-r^t}{1+r^t},
\tag{4}
\end{equation}
und dies ist, wie wir schon im \S~178 gesehen haben, eine Einheit des reellen Körpers $H_m$. Man erhält die conjugirten Werthe dieser Einheit in Bezug auf diesen Körper, wenn man $t\equiv5^\beta\pmod m$ setzt, und $\beta$ ein volles Restsystem nach dem Modul $\mu$ durchlaufen lässt. Wir setzen demnach:
\begin{equation}
\tau_\beta=\tan\frac{5^\beta\pi}{m}.
\tag{5}
\end{equation}
Nach \S~198, (7) ist $5^{\beta+\nu'}\equiv5^\beta+\mu\pmod m$, und daraus ergiebt sich nach einer bekannten trigonometrischen Formel:
\begin{equation}
\tau_{\beta+\nu'}=-\frac1{\tau_\beta}.
\tag{6}
\end{equation}
Setzen wir also nun
\begin{equation}
\lambda_\beta=\frac12\log\tau_\beta^2,
\tag{7}
\end{equation}
so ist nach (6)
\begin{equation}
\lambda_{\beta+\nu'}=-\lambda_\beta,
\tag{8}
\end{equation}
und demnach
\begin{equation}
\Theta^{\beta+\nu'}\lambda_{\beta+\nu'}=\Theta^\beta\lambda_\beta;
\tag{9}
\end{equation}
unter den in der Summe (3) vorkommenden Werthen von $\beta$ sind, wie wir schon im vorigen Paragraphen gesehen haben, nicht zwei nach dem Modul $\nu'$ congruent, und wenn wir daher mittelst
\srcpage{727}
der Formel (9) alle Werthe von $\beta$ auf das Intervall von 0 bis $\nu'-1$ reduciren, so bekommen wir alle Zahlen dieses Intervalles.

Demnach lässt sich die Formel (3) so darstellen:
\begin{equation}
X_1=-\frac2m(1,\Theta^{-1},r)\sum_{\beta=0}^{\nu'-1}\Theta^\beta\lambda_\beta.
\tag{10}
\end{equation}
Dies haben wir in dem im \S~197, (8) angegebenen Ausdrucke
\[
DB=2^{2^{\varkappa-4}(\varkappa-2)}\prod X_1=m^{\nu'}2^{-\nu'}\prod X_1
\]
einzusetzen, und dann ergiebt sich, wenn wir
\begin{equation}
U_\Theta=\sum_{\beta=0}^{\nu'-1}\Theta^\beta\lambda_\beta
\tag{11}
\end{equation}
setzen, mit Rücksicht auf \S~197, (14):
\begin{equation}
DB=\prod_\Theta U_\Theta.
\tag{12}
\end{equation}
Dies Product lässt sich nun auf folgende Weise durch eine Determinante darstellen, die nur die conjugirten Logarithmen der Einheiten $\tau$ enthält.

Lassen wir den Index $\Theta$ bei $U$ weg, so können wir das Gleichungssystem ansetzen:
\[
\begin{aligned}
U&=\lambda_0+\lambda_1\Theta+\lambda_2\Theta^2+\cdots+\lambda_{\nu'-1}\Theta^{\nu'-1},\\
\Theta U&=-\lambda_{\nu'-1}+\lambda_0\Theta+\lambda_1\Theta^2+\cdots+\lambda_{\nu'-2}\Theta^{\nu'-1},\\
&\hspace{2.5cm}\cdots\\
\Theta^{\nu'-1}U&=-\lambda_1-\lambda_2\Theta-\lambda_3\Theta^2+\cdots+\lambda_0\Theta^{\nu'-1},
\end{aligned}
\]
woraus durch Elimination der Potenzen von $\Theta$ folgt:
\begin{equation}
\begin{vmatrix}
\lambda_0-U&\lambda_1&\lambda_2&\cdots&\lambda_{\nu'-1}\\
-\lambda_{\nu'-1}&\lambda_0-U&\lambda_1&\cdots&\lambda_{\nu'-2}\\
\cdots&\cdots&\cdots&\cdots&\cdots\\
-\lambda_1&-\lambda_2&-\lambda_3&\cdots&\lambda_0-U
\end{vmatrix}=0.
\tag{13}
\end{equation}
Dies ist für $U$ eine Gleichung $\nu'$ten Grades, deren Wurzeln die Grössen $U_\Theta$ sind. Ordnet man nach Potenzen von $U$, so erhält, da $\nu'$ gerade ist, die höchste Potenz $U^{\nu'}$ von $U$ den Coefficienten 1, und man findet also das in (12) vorkommende Product der Wurzeln, wenn man $U=0$ setzt, aus der Determinante (13).

Die so gewonnene Determinante wollen wir etwas anders anordnen, indem wir die erste Zeile an ihrer Stelle lassen, die
\srcpage{728}
übrigen aber in umgekehrter Ordnung und mit veränderten Vorzeichen schreiben. Dabei sind $\frac12\nu'$ Vorzeichenänderungen nöthig, und wir erhalten daher für das Product (12) die durch folgende Gleichung definirte Determinante, deren absoluten Werth wir mit $T$ bezeichnen:
\begin{equation}
(-1)^{\frac12\nu'}
\begin{vmatrix}
\lambda_0&\lambda_1&\cdots&\lambda_{\nu'-2}&\lambda_{\nu'-1}\\
\lambda_1&\lambda_2&\cdots&\lambda_{\nu'-1}&-\lambda_0\\
\lambda_2&\lambda_3&\cdots&-\lambda_0&-\lambda_1\\
\cdots&\cdots&\cdots&\cdots&\cdots\\
\lambda_{\nu'-1}&-\lambda_0&\cdots&-\lambda_{\nu'-3}&-\lambda_{\nu'-2}
\end{vmatrix}=\pm T,
\tag{14}
\end{equation}
worin rechts von der von links unten nach rechts oben laufenden Diagonalreihe die negativen Zeichen stehen.

Die Determinante $T$ ist hier so geordnet, dass jede Zeile aus der vorangehenden durch die Substitution $(r,r^{5^\beta})$ entsteht. Jede Colonne enthält also ein System conjugirter Logarithmen.

Nach (12) ist jetzt
\begin{equation}
B=\frac{T}{D}.
\tag{15}
\end{equation}
Um über die Natur dieses Resultates Klarheit zu bekommen, ist es nöthig, die Zahl $D$ etwas genauer zu betrachten, die durch \S~195, (12), \S~197, (6) definirt war; und dies erfordert ein Eingehen auf die Theorie der Einheiten des Körpers $H_m$.

\end{document}
